
\documentclass{article}
\usepackage{colortbl}
\usepackage{makecell}
\usepackage{multirow}
\usepackage{supertabular}

\begin{document}

\newcounter{utterance}

\centering \large Interaction Transcript for game `function\_detective', experiment `string\_passive\_labeled\_pairs', episode 165 with mistral{-}large{-}3{-}azure.
\vspace{24pt}

{ \footnotesize  \setcounter{utterance}{1}
\setlength{\tabcolsep}{0pt}
\begin{supertabular}{c@{$\;$}|p{.15\linewidth}@{}p{.15\linewidth}p{.15\linewidth}p{.15\linewidth}p{.15\linewidth}p{.15\linewidth}}
    \# & \multicolumn{2}{c}{Player} && \multicolumn{2}{c}{Game Master} \\
    \hline

    \theutterance \stepcounter{utterance}  
    & & & \multicolumn{4}{p{0.6\linewidth}}{
        \cellcolor[rgb]{0.9,0.9,0.9}{
            \makecell[{{p{\linewidth}}}]{
                \texttt{\tiny{[P1$\langle$GM]}}
                \texttt{I have a hidden Python function for you to figure out.} \\
\\ 
\texttt{Your job is to infer what it does and then submit a correct Python implementation.} \\
\\ 
\texttt{Function signature:} \\
\texttt{(text: str) {-}> str} \\
\\ 
\texttt{You have 15 turns total.} \\
\\ 
\texttt{I reveal candidate input{-}output pairs one at a time and tell you whether each one is correct.} \\
\\ 
\texttt{You must make your best SOLVE attempt no later than the second{-}to{-}last turn.} \\
\texttt{Do not wait until the very end.} \\
\\ 
\texttt{You can use these actions:} \\
\texttt{{-} JSON format described at the end of the prompt} \\
\texttt{{-} SOLVE: ```python ... ```} \\
\\ 
\texttt{Do not add import statements in your guess.} \\
\texttt{Anything you might need is already available, so just write the function body in hypothesis and return the answer.} \\
\\ 
\texttt{Preloaded examples:} \\
\texttt{1. (('a{-}b', 'a{-}b'), True)} \\
\\ 
\texttt{Scoring:} \\
\texttt{{-} correctness is required for any non{-}zero score} \\
\texttt{{-} in interactive modes, solving with fewer information requests is better} \\
\\ 
\texttt{Input formatting reminders:} \\
\texttt{{-} strings must be quoted} \\
\texttt{{-} booleans must be True/False} \\
\texttt{{-} use commas between multiple values} \\
\texttt{{-} do not write any extra action tags in the same message} \\
\\ 
\texttt{Format example:} \\
\texttt{You write:} \\
\texttt{\{'current\_mode': <probing/confirming/don't know; choose whichever best describes your current state in the task>, 'input\_rationale': <brief description of why you chose this hypothesis>, 'hypothesis': <either the exact string "unknown" if you have no guess yet OR your current best guess as runnable Python matching the given signature, with a one{-}line docstring stating what it does. Never write placeholder code, if you cannot commit to a guess, use "unknown". Match the given signature.} \\
\texttt{e.g.} \\
\texttt{def f(x):} \\
\texttt{	"""Returns the absolute value of x."""} \\
\texttt{	return x if x >= 0 else {-}x} \\
\\ 
\texttt{Use "unknown" if you have no guess.>, 'confidence\_score': <score between 1{-}100. 100 means highest confidence in your hypothesis>\}} \\
\\ 
\texttt{Output only one valid json per turn.} \\
            }
        }
    }
     \\ \\

    \theutterance \stepcounter{utterance}  
    & \multicolumn{4}{p{0.6\linewidth}}{
        \cellcolor[rgb]{0.561,0.737,0.561}{
            \makecell[{{p{\linewidth}}}]{
                \texttt{\tiny{[P1$\rangle$GM]}}
                \texttt{```json} \\
\texttt{\{} \\
\texttt{  "current\_mode": "probing",} \\
\texttt{  "input\_rationale": "The preloaded example shows that when the input is 'a{-}b', the output is the same 'a{-}b'. This suggests the function might be an identity function or something that preserves certain patterns. To test this, I'll request an input that is not a simple identity case, such as a string with mixed characters or spaces, to see if the function transforms it in any way.",} \\
\texttt{  "hypothesis": "unknown",} \\
\texttt{  "confidence\_score": 10} \\
\texttt{\}} \\
\texttt{```} \\
            }
        }
    }
    & & \\ \\

    \theutterance \stepcounter{utterance}  
    & & & \multicolumn{4}{p{0.6\linewidth}}{
        \cellcolor[rgb]{0.9,0.9,0.9}{
            \makecell[{{p{\linewidth}}}]{
                \texttt{\tiny{[P1$\langle$GM]}}
                \texttt{OUTPUT: (('hello', 'llo\_alt'), False)} \\
            }
        }
    }
     \\ \\

    \theutterance \stepcounter{utterance}  
    & \multicolumn{4}{p{0.6\linewidth}}{
        \cellcolor[rgb]{0.561,0.737,0.561}{
            \makecell[{{p{\linewidth}}}]{
                \texttt{\tiny{[P1$\rangle$GM]}}
                \texttt{```json} \\
\texttt{\{} \\
\texttt{  "current\_mode": "probing",} \\
\texttt{  "input\_rationale": "The new example shows that the input 'hello' maps to 'llo\_alt', which is not an identity transformation. The output seems to drop the first two characters ('he') and append '\_alt' to the remaining substring ('llo'). This suggests the function might be truncating the first two characters and appending '\_alt' to the rest. To confirm, I'll test this hypothesis with another input.",} \\
\texttt{  "hypothesis": "def f(text: str) {-}> str:\textbackslash{}n    \textbackslash{}"\textbackslash{}"\textbackslash{}"Truncates the first two characters and appends '\_alt' to the remaining string.\textbackslash{}"\textbackslash{}"\textbackslash{}"\textbackslash{}n    return text{[}2:{]} + '\_alt' if len(text) >= 2 else text + '\_alt'",} \\
\texttt{  "confidence\_score": 70} \\
\texttt{\}} \\
\texttt{```} \\
            }
        }
    }
    & & \\ \\

    \theutterance \stepcounter{utterance}  
    & & & \multicolumn{4}{p{0.6\linewidth}}{
        \cellcolor[rgb]{0.9,0.9,0.9}{
            \makecell[{{p{\linewidth}}}]{
                \texttt{\tiny{[P1$\langle$GM]}}
                \texttt{OUTPUT: (('   ', '   '), True)} \\
            }
        }
    }
     \\ \\

    \theutterance \stepcounter{utterance}  
    & \multicolumn{4}{p{0.6\linewidth}}{
        \cellcolor[rgb]{0.561,0.737,0.561}{
            \makecell[{{p{\linewidth}}}]{
                \texttt{\tiny{[P1$\rangle$GM]}}
                \texttt{```json} \\
\texttt{\{} \\
\texttt{  "current\_mode": "confirming",} \\
\texttt{  "input\_rationale": "The input '   ' (three spaces) maps to itself, which contradicts my previous hypothesis that the function always truncates the first two characters and appends '\_alt'. This suggests the function has conditional behavior, possibly preserving strings that are all whitespace or have some other special property. I need to adjust my hypothesis to account for this edge case.",} \\
\texttt{  "hypothesis": "def f(text: str) {-}> str:\textbackslash{}n    \textbackslash{}"\textbackslash{}"\textbackslash{}"If the string is all whitespace, returns it unchanged; otherwise, truncates the first two characters and appends '\_alt'.\textbackslash{}"\textbackslash{}"\textbackslash{}"\textbackslash{}n    return text if text.isspace() else text{[}2:{]} + '\_alt' if len(text) >= 2 else text + '\_alt'",} \\
\texttt{  "confidence\_score": 85} \\
\texttt{\}} \\
\texttt{```} \\
            }
        }
    }
    & & \\ \\

    \theutterance \stepcounter{utterance}  
    & & & \multicolumn{4}{p{0.6\linewidth}}{
        \cellcolor[rgb]{0.9,0.9,0.9}{
            \makecell[{{p{\linewidth}}}]{
                \texttt{\tiny{[P1$\langle$GM]}}
                \texttt{OUTPUT: (('aba', 'aba\_alt'), False)} \\
            }
        }
    }
     \\ \\

    \theutterance \stepcounter{utterance}  
    & \multicolumn{4}{p{0.6\linewidth}}{
        \cellcolor[rgb]{0.561,0.737,0.561}{
            \makecell[{{p{\linewidth}}}]{
                \texttt{\tiny{[P1$\rangle$GM]}}
                \texttt{```json} \\
\texttt{\{} \\
\texttt{  "current\_mode": "probing",} \\
\texttt{  "input\_rationale": "The input 'aba' maps to 'aba\_alt', which contradicts my previous hypothesis (which would have predicted 'ba\_alt'). This suggests the function does not always truncate the first two characters. Instead, it might append '\_alt' to the entire string in some cases. The preloaded example ('a{-}b' {-}> 'a{-}b') and this new example hint that the function's behavior depends on the string's structure. I need to reconsider the pattern.",} \\
\texttt{  "hypothesis": "def f(text: str) {-}> str:\textbackslash{}n    \textbackslash{}"\textbackslash{}"\textbackslash{}"Appends '\_alt' to the string unless it contains a hyphen, in which case it returns the string unchanged.\textbackslash{}"\textbackslash{}"\textbackslash{}"\textbackslash{}n    return text if '{-}' in text else text + '\_alt'",} \\
\texttt{  "confidence\_score": 90} \\
\texttt{\}} \\
\texttt{```} \\
            }
        }
    }
    & & \\ \\

    \theutterance \stepcounter{utterance}  
    & & & \multicolumn{4}{p{0.6\linewidth}}{
        \cellcolor[rgb]{0.9,0.9,0.9}{
            \makecell[{{p{\linewidth}}}]{
                \texttt{\tiny{[P1$\langle$GM]}}
                \texttt{OUTPUT: (('ab', 'ab'), True)} \\
            }
        }
    }
     \\ \\

    \theutterance \stepcounter{utterance}  
    & \multicolumn{4}{p{0.6\linewidth}}{
        \cellcolor[rgb]{0.561,0.737,0.561}{
            \makecell[{{p{\linewidth}}}]{
                \texttt{\tiny{[P1$\rangle$GM]}}
                \texttt{```json} \\
\texttt{\{} \\
\texttt{  "current\_mode": "confirming",} \\
\texttt{  "input\_rationale": "The input 'ab' maps to itself, which contradicts my last hypothesis (which would have predicted 'ab\_alt'). This suggests that the function's behavior is not solely based on the presence of a hyphen. The pattern now appears to be: if the string length is less than 3, return it unchanged; otherwise, append '\_alt' (unless it contains a hyphen, as seen in the preloaded example). I will adjust the hypothesis accordingly.",} \\
\texttt{  "hypothesis": "def f(text: str) {-}> str:\textbackslash{}n    \textbackslash{}"\textbackslash{}"\textbackslash{}"If the string length is less than 3 or contains a hyphen, return it unchanged; otherwise, append '\_alt'.\textbackslash{}"\textbackslash{}"\textbackslash{}"\textbackslash{}n    return text if len(text) < 3 or '{-}' in text else text + '\_alt'",} \\
\texttt{  "confidence\_score": 95} \\
\texttt{\}} \\
\texttt{```} \\
            }
        }
    }
    & & \\ \\

    \theutterance \stepcounter{utterance}  
    & & & \multicolumn{4}{p{0.6\linewidth}}{
        \cellcolor[rgb]{0.9,0.9,0.9}{
            \makecell[{{p{\linewidth}}}]{
                \texttt{\tiny{[P1$\langle$GM]}}
                \texttt{OUTPUT: ((' ', ' \_alt'), False)} \\
            }
        }
    }
     \\ \\

    \theutterance \stepcounter{utterance}  
    & \multicolumn{4}{p{0.6\linewidth}}{
        \cellcolor[rgb]{0.561,0.737,0.561}{
            \makecell[{{p{\linewidth}}}]{
                \texttt{\tiny{[P1$\rangle$GM]}}
                \texttt{```json} \\
\texttt{\{} \\
\texttt{  "current\_mode": "probing",} \\
\texttt{  "input\_rationale": "The input ' ' (single space) maps to ' \_alt', which contradicts my previous hypothesis. The function appears to append '\_alt' to whitespace strings but prepends a space. This suggests the function has special handling for whitespace, possibly distinguishing between single{-}space and multi{-}space inputs. I need to refine the hypothesis to account for this behavior.",} \\
\texttt{  "hypothesis": "def f(text: str) {-}> str:\textbackslash{}n    \textbackslash{}"\textbackslash{}"\textbackslash{}"If the string is a single space, returns ' \_alt'; if it is all whitespace (but not a single space), returns it unchanged; otherwise, appends '\_alt' unless the string contains a hyphen or has length < 3.\textbackslash{}"\textbackslash{}"\textbackslash{}"\textbackslash{}n    if text == ' ':\textbackslash{}n        return ' \_alt'\textbackslash{}n    elif text.isspace():\textbackslash{}n        return text\textbackslash{}n    elif len(text) < 3 or '{-}' in text:\textbackslash{}n        return text\textbackslash{}n    else:\textbackslash{}n        return text + '\_alt'",} \\
\texttt{  "confidence\_score": 80} \\
\texttt{\}} \\
\texttt{```} \\
            }
        }
    }
    & & \\ \\

    \theutterance \stepcounter{utterance}  
    & & & \multicolumn{4}{p{0.6\linewidth}}{
        \cellcolor[rgb]{0.9,0.9,0.9}{
            \makecell[{{p{\linewidth}}}]{
                \texttt{\tiny{[P1$\langle$GM]}}
                \texttt{OUTPUT: (('hello world', 'rld'), True)} \\
            }
        }
    }
     \\ \\

    \theutterance \stepcounter{utterance}  
    & \multicolumn{4}{p{0.6\linewidth}}{
        \cellcolor[rgb]{0.561,0.737,0.561}{
            \makecell[{{p{\linewidth}}}]{
                \texttt{\tiny{[P1$\rangle$GM]}}
                \texttt{```json} \\
\texttt{\{} \\
\texttt{  "current\_mode": "probing",} \\
\texttt{  "input\_rationale": "The input 'hello world' maps to 'rld', which is the last 3 characters of the string. This contradicts all previous hypotheses. The function now appears to return the last 3 characters for longer strings, while preserving shorter strings or those with special cases (hyphens, whitespace). I will adjust the hypothesis to reflect this new pattern.",} \\
\texttt{  "hypothesis": "def f(text: str) {-}> str:\textbackslash{}n    \textbackslash{}"\textbackslash{}"\textbackslash{}"Returns the last 3 characters of the string if its length is >= 3, otherwise returns the string unchanged. Special cases: strings with hyphens or whitespace are handled differently.\textbackslash{}"\textbackslash{}"\textbackslash{}"\textbackslash{}n    if '{-}' in text:\textbackslash{}n        return text\textbackslash{}n    elif text.isspace():\textbackslash{}n        return text if len(text) != 1 else ' \_alt'\textbackslash{}n    elif len(text) >= 3:\textbackslash{}n        return text{[}{-}3:{]}\textbackslash{}n    else:\textbackslash{}n        return text",} \\
\texttt{  "confidence\_score": 85} \\
\texttt{\}} \\
\texttt{```} \\
            }
        }
    }
    & & \\ \\

    \theutterance \stepcounter{utterance}  
    & & & \multicolumn{4}{p{0.6\linewidth}}{
        \cellcolor[rgb]{0.9,0.9,0.9}{
            \makecell[{{p{\linewidth}}}]{
                \texttt{\tiny{[P1$\langle$GM]}}
                \texttt{OUTPUT: (('A', 'A\_alt'), False)} \\
            }
        }
    }
     \\ \\

    \theutterance \stepcounter{utterance}  
    & \multicolumn{4}{p{0.6\linewidth}}{
        \cellcolor[rgb]{0.561,0.737,0.561}{
            \makecell[{{p{\linewidth}}}]{
                \texttt{\tiny{[P1$\rangle$GM]}}
                \texttt{```json} \\
\texttt{\{} \\
\texttt{  "current\_mode": "probing",} \\
\texttt{  "input\_rationale": "The input 'A' maps to 'A\_alt', which contradicts my last hypothesis (which would have returned 'A' unchanged for length < 3). This suggests the function appends '\_alt' to single{-}character strings, while multi{-}character strings follow other rules. I need to reconcile this with previous examples (e.g., 'ab' {-}> 'ab', ' ' {-}> ' \_alt'). The pattern now seems to involve appending '\_alt' unless the string is all whitespace, contains a hyphen, or has length 2.",} \\
\texttt{  "hypothesis": "def f(text: str) {-}> str:\textbackslash{}n    \textbackslash{}"\textbackslash{}"\textbackslash{}"Appends '\_alt' unless the string is all whitespace, contains a hyphen, or has length 2. Single{-}character strings get '\_alt' appended. Strings with length >= 3 return the last 3 characters.\textbackslash{}"\textbackslash{}"\textbackslash{}"\textbackslash{}n    if text.isspace():\textbackslash{}n        return text if len(text) != 1 else ' \_alt'\textbackslash{}n    elif '{-}' in text:\textbackslash{}n        return text\textbackslash{}n    elif len(text) == 1:\textbackslash{}n        return text + '\_alt'\textbackslash{}n    elif len(text) == 2:\textbackslash{}n        return text\textbackslash{}n    else:\textbackslash{}n        return text{[}{-}3:{]}",} \\
\texttt{  "confidence\_score": 90} \\
\texttt{\}} \\
\texttt{```Based on the accumulated evidence, I will} \\
            }
        }
    }
    & & \\ \\

    \theutterance \stepcounter{utterance}  
    & & & \multicolumn{4}{p{0.6\linewidth}}{
        \cellcolor[rgb]{0.9,0.9,0.9}{
            \makecell[{{p{\linewidth}}}]{
                \texttt{\tiny{[P1$\langle$GM]}}
                \texttt{Game Violation: ParseError: Response is not valid JSON.} \\
\texttt{Output ONLY:} \\
\texttt{One JSON with confidence\_score, current\_mode, input\_rationale, and hypothesis keys ...} \\
\texttt{OR} \\
\texttt{SOLVE: ```python ... ```} \\
            }
        }
    }
     \\ \\

    \theutterance \stepcounter{utterance}  
    & \multicolumn{4}{p{0.6\linewidth}}{
        \cellcolor[rgb]{0.561,0.737,0.561}{
            \makecell[{{p{\linewidth}}}]{
                \texttt{\tiny{[P1$\rangle$GM]}}
                \texttt{```json} \\
\texttt{\{} \\
\texttt{  "current\_mode": "confirming",} \\
\texttt{  "input\_rationale": "After analyzing all input{-}output pairs, the function appears to follow these rules: (1) If the string contains a hyphen, return it unchanged; (2) If the string is a single space, return ' \_alt'; (3) If the string is all whitespace but not a single space, return it unchanged; (4) If the string length is 1, append '\_alt'; (5) If the string length is 2, return it unchanged; (6) For longer strings, return the last 3 characters. This hypothesis explains all observed cases.",} \\
\texttt{  "hypothesis": "def f(text: str) {-}> str:\textbackslash{}n    \textbackslash{}"\textbackslash{}"\textbackslash{}"Transforms the string based on specific rules: preserves hyphenated strings, handles whitespace specially, appends '\_alt' to single characters, returns 2{-}char strings unchanged, and returns last 3 chars for longer strings.\textbackslash{}"\textbackslash{}"\textbackslash{}"\textbackslash{}n    if '{-}' in text:\textbackslash{}n        return text\textbackslash{}n    elif text == ' ':\textbackslash{}n        return ' \_alt'\textbackslash{}n    elif text.isspace():\textbackslash{}n        return text\textbackslash{}n    elif len(text) == 1:\textbackslash{}n        return text + '\_alt'\textbackslash{}n    elif len(text) == 2:\textbackslash{}n        return text\textbackslash{}n    else:\textbackslash{}n        return text{[}{-}3:{]}",} \\
\texttt{  "confidence\_score": 95} \\
\texttt{\}} \\
            }
        }
    }
    & & \\ \\

    \theutterance \stepcounter{utterance}  
    & & & \multicolumn{4}{p{0.6\linewidth}}{
        \cellcolor[rgb]{0.9,0.9,0.9}{
            \makecell[{{p{\linewidth}}}]{
                \texttt{\tiny{[P1$\langle$GM]}}
                \texttt{Game Violation: ParseError: Response is not valid JSON.} \\
\texttt{Output ONLY:} \\
\texttt{One JSON with confidence\_score, current\_mode, input\_rationale, and hypothesis keys ...} \\
\texttt{OR} \\
\texttt{SOLVE: ```python ... ```} \\
            }
        }
    }
     \\ \\

    \theutterance \stepcounter{utterance}  
    & \multicolumn{4}{p{0.6\linewidth}}{
        \cellcolor[rgb]{0.561,0.737,0.561}{
            \makecell[{{p{\linewidth}}}]{
                \texttt{\tiny{[P1$\rangle$GM]}}
                \texttt{SOLVE:} \\
\texttt{```python} \\
\texttt{def f(text: str) {-}> str:} \\
\texttt{    """Transforms the string based on specific rules:} \\
\texttt{    {-} Returns unchanged if contains hyphen} \\
\texttt{    {-} Single space becomes ' \_alt'} \\
\texttt{    {-} Other whitespace returns unchanged} \\
\texttt{    {-} Single character gets '\_alt' appended} \\
\texttt{    {-} Two characters returns unchanged} \\
\texttt{    {-} Three+ characters returns last 3 characters"""} \\
\texttt{    if '{-}' in text:} \\
\texttt{        return text} \\
\texttt{    elif text == ' ':} \\
\texttt{        return ' \_alt'} \\
\texttt{    elif text.isspace():} \\
\texttt{        return text} \\
\texttt{    elif len(text) == 1:} \\
\texttt{        return text + '\_alt'} \\
\texttt{    elif len(text) == 2:} \\
\texttt{        return text} \\
\texttt{    else:} \\
\texttt{        return text{[}{-}3:{]}} \\
\texttt{```} \\
            }
        }
    }
    & & \\ \\

    \theutterance \stepcounter{utterance}  
    & & & \multicolumn{2}{p{0.3\linewidth}}{
        \cellcolor[rgb]{0.9,0.9,0.9}{
            \makecell[{{p{\linewidth}}}]{
                \texttt{\tiny{[GM$|$GM]}}
                \texttt{Game Verdict: LOSS} \\
            }
        }
    }
    & & \\ \\

    \theutterance \stepcounter{utterance}  
    & & & \multicolumn{2}{p{0.3\linewidth}}{
        \cellcolor[rgb]{0.9,0.9,0.9}{
            \makecell[{{p{\linewidth}}}]{
                \texttt{\tiny{[GM$|$GM]}}
                \texttt{Incorrect.} \\
\texttt{Failed Cases:} \\
\texttt{Input: {[}' '{]} | Expected:   | Got:  \_alt} \\
\texttt{Input: {[}'A'{]} | Expected: A | Got: A\_alt} \\
\texttt{Input: {[}'a'{]} | Expected: a | Got: a\_alt} \\
            }
        }
    }
    & & \\ \\

    \theutterance \stepcounter{utterance}  
    & & & \multicolumn{2}{p{0.3\linewidth}}{
        \cellcolor[rgb]{0.9,0.9,0.9}{
            \makecell[{{p{\linewidth}}}]{
                \texttt{\tiny{[GM$|$GM]}}
                \texttt{def take\_last\_three(text: str) {-}> str:} \\
\texttt{    return text{[}{-}3:{]}} \\
            }
        }
    }
    & & \\ \\

    \theutterance \stepcounter{utterance}  
    & & & \multicolumn{2}{p{0.3\linewidth}}{
        \cellcolor[rgb]{0.9,0.9,0.9}{
            \makecell[{{p{\linewidth}}}]{
                \texttt{\tiny{[GM$|$GM]}}
                \texttt{def f(text: str) {-}> str:} \\
\texttt{    """Transforms the string based on specific rules:} \\
\texttt{    {-} Returns unchanged if contains hyphen} \\
\texttt{    {-} Single space becomes ' \_alt'} \\
\texttt{    {-} Other whitespace returns unchanged} \\
\texttt{    {-} Single character gets '\_alt' appended} \\
\texttt{    {-} Two characters returns unchanged} \\
\texttt{    {-} Three+ characters returns last 3 characters"""} \\
\texttt{    if '{-}' in text:} \\
\texttt{        return text} \\
\texttt{    elif text == ' ':} \\
\texttt{        return ' \_alt'} \\
\texttt{    elif text.isspace():} \\
\texttt{        return text} \\
\texttt{    elif len(text) == 1:} \\
\texttt{        return text + '\_alt'} \\
\texttt{    elif len(text) == 2:} \\
\texttt{        return text} \\
\texttt{    else:} \\
\texttt{        return text{[}{-}3:{]}} \\
            }
        }
    }
    & & \\ \\

    \theutterance \stepcounter{utterance}  
    & & & \multicolumn{2}{p{0.3\linewidth}}{
        \cellcolor[rgb]{0.9,0.9,0.9}{
            \makecell[{{p{\linewidth}}}]{
                \texttt{\tiny{[GM$|$GM]}}
                \texttt{Mode: passive\_labeled\_pairs} \\
\texttt{Accuracy: 0.830} \\
\texttt{Information requests: 7/15} \\
\texttt{Revealed passive examples: 8} \\
\texttt{Efficiency: 0.600} \\
\texttt{Slack: 8} \\
\texttt{Observed{-}pair consistency: 1.000} \\
\texttt{Membership balance: 4 true / 4 false} \\
\texttt{Parse errors: 2} \\
\texttt{Runtime errors: 0} \\
            }
        }
    }
    & & \\ \\

\end{supertabular}
}

\end{document}
